\section{Experiments}
\label{sec:experiments}

All experiments use a decoder-only language model at GPT-2 scale unless
otherwise noted.
Evaluation metric: exact-match trace accuracy on a held-out set of 1\,000
programs per difficulty level, plus backtrack accuracy (fraction of
\textsc{Fork} nodes where the correct branch is selected).

\subsection{E1: SFT Baseline}
\label{sec:e1}

\paragraph{Setup.}
We fine-tune on \gore{} curriculum data (levels 0--4, ${\approx}$100K programs
total) using maximum-likelihood supervision on full execution traces.
Training proceeds for 3 epochs with a cosine learning-rate schedule; we use the
trace format described in Section~\ref{sec:language}.

\paragraph{Metrics.}
Exact-match trace accuracy; per-level accuracy breakdown;
backtrack accuracy; trace length distribution.

\paragraph{Expected outcome.}
SFT should achieve high accuracy at levels 0--2 and degrade at levels 3--4,
motivating the RL stage.

\todo{Fill in actual numbers; add learning curves; ablation over curriculum
      ordering (random vs.\ staged).}

\subsection{E2: Dense vs.\ Sparse Reward (Dr.\,GRPO)}
\label{sec:e2}

We study three conditions:

\paragraph{E2a: GRPO + sparse reward.}
Standard GRPO \citep{shao2024deepseekmath} with binary reward: $+1$ if the
final answer is correct, $0$ otherwise.
This is the baseline RL condition.

\paragraph{E2b: Dr.\,GRPO + dense reward.}
Dr.\,GRPO with the per-step interpreter reward from
Section~\ref{sec:interpreter-prm}.
This is the main experimental condition.
The comparison E2a vs.\ E2b isolates the value of dense reward; the comparison
E1 vs.\ E2b isolates the value of RL over SFT.

\paragraph{E2c: DAPO.}
DAPO \citep{yu2025dapo} as a memory-efficient ablation (no KL penalty, dynamic
clip, token-level policy gradient).
Same sparse reward as E2a to control for the reward signal.

\paragraph{Metrics.}
Trace accuracy per level; backtrack accuracy; reward curve; KL divergence from
SFT initialisation; GPU memory consumption.

\todo{Fill in results; add per-level accuracy table; add qualitative examples
      of emergent backtracking behaviour.}

\subsection{E3: Spectral Analysis (Spectral-R1)}
\label{sec:e3}

We apply PCA and SVD to the residual-stream activations of the Dr.\,GRPO model,
conditioned on the node type of the \emph{current} trace token
(\textsc{Fork}, \textsc{Cut}, or \textsc{Step}).

\paragraph{Setup.}
We collect activations at the final residual stream position for 10\,000 trace
tokens (balanced across node types) from the E2b model.
We fit PCA jointly on all three conditions and inspect the leading principal
components.

\paragraph{Intervention.}
We ablate the direction most predictive of \textsc{Fork} vs.\ \textsc{Step}
(identified by a linear probe) and measure the drop in backtrack accuracy.

\paragraph{Expected outcome.}
We expect \textsc{Fork}/\textsc{Cut} activations to occupy a distinguishable
subspace from \textsc{Step} activations, and ablation of this subspace to cause
a selective drop in backtrack accuracy without affecting \textsc{Step} accuracy.

\todo{Add PCA plots; confusion matrix of linear probe; ablation table.}

\subsection{E4: Transfer to Natural-Language Reasoning}
\label{sec:e4}

We evaluate the E2b model (and an SFT-only baseline) on natural-language
reasoning benchmarks after adapter fine-tuning on a small NL dataset.

\paragraph{Benchmarks.}
GSM8K (arithmetic multi-step), ARC-Challenge (science QA), and a
backtracking-heavy subset of MATH.

\paragraph{Setup.}
Zero-shot and 5-shot evaluation; comparison to an equivalent-compute baseline
trained on NL data only.

\paragraph{Expected outcome.}
\gore{}-pretrained models should show improved performance on tasks that require
systematic search (multi-step arithmetic with correction, proof-by-cases), with
smaller gains on tasks that are primarily forward-only.

\todo{Fill in benchmark numbers; add breakdown by task type; discuss
      failure cases.}
